\documentclass[11pt,a4paper]{article}
\usepackage[utf8]{inputenc}
\usepackage[T1]{fontenc}
\usepackage{lmodern}
\usepackage[margin=2.6cm]{geometry}
\usepackage{amsmath,amssymb}
\usepackage{booktabs}
\usepackage{array}
\usepackage{enumitem}
\usepackage[colorlinks=true,linkcolor=blue,urlcolor=blue]{hyperref}
\setlength{\parindent}{0pt}
\setlength{\parskip}{6pt}

\newcommand{\HS}{\mathrm{HS}}
\newcommand{\LS}{\mathrm{LS}}
\newcommand{\Dq}{Dq}
\newcommand{\kb}{k_{\mathrm{B}}}

\title{\bfseries Guía para arreglar \texttt{sco0} y crear la base para continuar}
\author{Alejandro Rodríguez}
\date{Noviembre 2026}

\begin{document}
\maketitle

\begin{quote}
\emph{Evalúa lo que generó Gemini (\texttt{sco0/}), dice qué conservar y qué
reemplazar, y especifica el modelo objetivo para retomar el doctorado.
Complementa a la ``Memoria del proyecto original''.}
\end{quote}

\tableofcontents

\section{Diagnóstico de \texttt{sco0}}

\subsection{Lo que está bien (conservar)}

\begin{itemize}
\item \textbf{Estructura limpia y modular}: \texttt{config.py}, \texttt{model.py},
  \texttt{equilibrium.py}, \texttt{monte\_carlo.py}, \texttt{protocol.py},
  \texttt{output.py}.
\item \textbf{Reproducibilidad}: semilla fija, \texttt{numpy.random.SeedSequence}.
\item \textbf{Configuración declarativa} en JSON con validación estricta.
\item \textbf{Referencia cuasiestática} con \texttt{scipy.optimize}.
\item \textbf{Tests} y \textbf{salidas CSV/PNG} automatizadas.
\item \textbf{Honestidad documental}: deja claro qué \emph{no} modela.
\end{itemize}

\subsection{Lo que perdió (el problema real)}

\begin{center}
\begin{tabular}{p{4.6cm}p{4.6cm}p{4.6cm}}
\toprule
Aspecto & Tu modelo original & \texttt{sco0} (Gemini) \\
\midrule
Temperatura & Metropolis a $T$ finita (Boltzmann) & 0 K (descenso cuesta abajo) \\
Espín & Se propone y se acepta con $e^{-\Delta E/\kb T}$ & Nunca se propone; se recalcula como suelo \\
N.\ de centros & 1 (\texttt{prueba}) y red $N^3$ (\texttt{project\_sco}) & 1--2 \\
Cooperatividad & Red elástica + intercambio ($j_{\mathrm{isin}}$) & Solo puente $K_{\mathrm{bridge}}$ entre 2 Fe \\
Energía electrónica & \textbf{Ligand-field $\Dq$} (Tanabe--Sugano linealizado) & ``Potencial de Jørgensen'' ad hoc \\
Histéresis & Objetivo (ciclo térmico) & Bucle a 0 K = atrapamiento \\
\bottomrule
\end{tabular}
\end{center}

\textbf{Conclusión:} Gemini escribió un código limpio pero \emph{no es tu
modelo}: es una optimización de geometría a 0 K con un potencial inventado. Hay
que \textbf{conservar la estructura} y \textbf{reemplazar la física}.

\subsection{Lo que está mal en \texttt{sco0}}

\begin{enumerate}
\item \textbf{0 K como ``Monte Carlo''.} \texttt{monte\_carlo.py} acepta solo
  $E_{\mathrm{trial}}\le E_{\mathrm{old}}$; no hay Boltzmann.
\item \textbf{El espín es esclavo de la geometría.} \texttt{ground\_state()} elige
  $\HS$/$\LS$ por mínimo a geometría fija; no hay variable de espín muestreada.
\item \textbf{``Temperatura nominal'' sin sentido físico.} $T$ es una etiqueta,
  nunca entra en $\kb T$.
\item \textbf{El ``potencial de Jørgensen'' no es tu modelo.}
  $E = E_{\mathrm{scale}}\bigl(\tfrac{1}{9x^9} - 0.02\,\rho\,\ell/x^5 - 1/x\bigr)+P_s$
  es una Lennard-Jones modificada, no el cruce $-4\Dq$ vs $-24\Dq+P$.
\item \textbf{Unidades poco físicas} ($P_{\LS}=3.9$ eV, $E_{\mathrm{scale}}=15$ eV).
\end{enumerate}

\section{El modelo objetivo (la base a construir)}

Unifica las tres versiones previas. La idea central (de la tesis y de
\texttt{prueba/main.cpp}):

\begin{quote}
\emph{La temperatura contrae la red $\to$ sube el campo de ligandos $\Dq$ $\to$
cruza $\HS$/$\LS$.}
\end{quote}

\subsection{Energía de un centro (campo de ligandos, versión corregida)}

Sea $r$ la distancia Fe--N. El campo de ligandos decae con la distancia (modelo
de cargas puntuales, $10\Dq\propto r^{-5}$):
\begin{equation}
\Dq(r) = \Dq_0\,\bigl(r_0/r\bigr)^5 .
\end{equation}

Energía electrónica de cada estado (linealización de Tanabe--Sugano):
\begin{align}
E_{\mathrm{el}}(r,\HS) &= -4\,\Dq(r)\\
E_{\mathrm{el}}(r,\LS) &= -24\,\Dq(r) + P
\end{align}

El cruce está en $\Dq(r_c)=P/20$, es decir $r_c = r_0\,(20\Dq_0/P)^{1/5}$.
Eligiendo $P$ para que $r_c$ quede entre $r_{\HS}$ y $r_{\LS}$ se obtiene enlace
\emph{largo en $\HS$} y \emph{corto en $\LS$} (la contracción de red). Unidades:
cm$^{-1}$ + K con $\kb=0.695$ cm$^{-1}$/K (o eV con $\kb=8.617\times10^{-5}$ eV/K).
Valores de referencia (Fe$^{2+}$): $\Dq_0\approx 1200$--$2000$ cm$^{-1}$ en
$r_0\approx 2.0$ Å, $B\approx 917$--$1050$ cm$^{-1}$, $C\approx 4B$,
$P\approx 20000$--$25000$ cm$^{-1}$ (ajustable).

\subsection{Acoplamiento elástico (cooperatividad)}

Para $N$ centros, resorte entre vecinos (recupera \texttt{project\_sco}):
\begin{equation}
E_{\mathrm{elast}} = \frac{K}{2}\sum_{(i,j)}\bigl(r_{ij}-(r_i+r_j)\bigr)^2 ,
\end{equation}
o, para empezar, un acoplamiento al promedio de la red (campo medio elástico):
\begin{equation}
E_{\mathrm{lat}} = \frac{K_{\mathrm{red}}}{2}\sum_i\bigl(r_i - R_{\mathrm{lat}}\bigr)^2 .
\end{equation}
El término de intercambio $-J\sum_{(i,j)}\sigma_i\sigma_j$ (opcional) añade otra
vía de cooperatividad.

\subsection{Metropolis a temperatura finita (lo que falta)}

Por cada paso MC, \textbf{proponer} (i) voltear el espín de un centro y/o (ii)
mover una distancia, calcular $\Delta E$ \emph{una sola vez}, y aceptar con
\begin{equation}
p = \min\!\bigl(1,\ e^{-\Delta E/\kb T}\bigr) .
\end{equation}

Solo así se recupera el ensemble térmico y, con un ciclo $T$, la
\textbf{histéresis} cooperativa.

\subsection{Protocolo (ciclo térmico real)}

Barrido de temperatura $T$ (no ``nominal''), con la red libre (NPT) o fija (NVT),
en dos direcciones $\Rightarrow$ fracción $\HS$ vs $T$ con lazo.

\section{Hoja de ruta en fases}

\noindent\textbf{Estado actual:} las Fases 0 y 1 ya están implementadas en el
paquete \texttt{sco0/} (campo de ligandos $\Dq(r)$, constantes físicas,
\texttt{metropolis.py} a temperatura finita, protocolos de distancia y de
temperatura, pruebas y figuras). Las Fases 2--4 son el trabajo por delante.

\begin{description}
\item[Fase 0 -- Unidades y constantes.] \textbf{[hecho]} \texttt{sco0/constants.py}
  con $\kb$, $B$, $C$, $Dq_0$; unidad única cm$^{-1}$ + K.
\item[Fase 1 -- Un centro, $T$ finita.] \textbf{[hecho]}
  \texttt{model.py} con la energía de campo de ligandos, \texttt{metropolis.py}
  con propuesta de espín + Boltzmann, y los protocolos \texttt{distance} y
  \texttt{temperature}. Objetivo logrado: cruce suave validado contra la solución
  analítica (dos niveles).
\item[Fase 2 -- Cooperatividad.] $N$ centros (1D $\to$ 2D/3D), resorte entre
  vecinos, Metropolis conjunto (espín + posición + volumen). Objetivo: transición
  de primer orden y lazo de histéresis.
\item[Fase 3 -- Ciclo térmico y validación.] Calentar/enfriar, fracción $\HS$ vs
  $T$, ancho de histéresis; comparar con los datos experimentales
  (\texttt{~/Dropbox/PhD/2018/Semestre-I/SCO\_DATA/}).
\item[Fase 4 (opcional) -- Refinamiento.] Tabla completa de Tanabe--Sugano
  interpolada en $\Dq/B$, o parámetros de DFT (Paulsen--Trautwein 2004).
\end{description}

\section{Qué tocar en el código (mapa concreto)}

La columna ``Acción'' de las Fases 0--1 ya está aplicada. El paquete actual
quedó así (todo lo marcado como \textbf{hecho} está implementado):

\begin{center}
\begin{tabular}{ll}
\toprule
Archivo & Estado \\
\midrule
\texttt{sco0/constants.py} & \textbf{hecho}: $\kb$, Racah $B$, $C$ \\
\texttt{sco0/config.py} & \textbf{hecho}: $\Dq_0$, $r_0$, $P$, protocolos, MC, semilla \\
\texttt{sco0/model.py} & \textbf{hecho}: $\Dq(r)$, $E_{\HS}$, $E_{\LS}$, cruce, fracción analítica \\
\texttt{sco0/metropolis.py} & \textbf{hecho}: espín propuesto, $e^{-\Delta E/\kb T}$ \\
\texttt{sco0/monte\_carlo.py} & \textbf{hecho}: Metropolis térmico con réplicas \\
\texttt{sco0/equilibrium.py} & \textbf{hecho}: referencia analítica de dos niveles \\
\texttt{sco0/protocol.py} & \textbf{hecho}: \texttt{distance} y \texttt{temperature} \\
\texttt{sco0/output.py} & \textbf{hecho}: \texttt{surfaces.png} y \texttt{crossover.png} \\
\texttt{tests/} & \textbf{hecho}: 14 pruebas (cruce, Boltzmann, reproducción) \\
\midrule
\textbf{Pendiente (Fase 2)} & $N$ centros, resorte entre vecinos, acoplamiento $J$ \\
\bottomrule
\end{tabular}
\end{center}

\section{Criterios de ``hecho'' (verificables)}

\begin{enumerate}
\item \texttt{python -m unittest discover -s tests -v} en verde.
\item Un centro a $T$ finita da fracción $\HS$ \textbf{suave} en $T$.
\item $N$ centros muestran \textbf{histéresis} (lazo en el ciclo), no atrapamiento.
\item Parámetros ($\Dq$, $B$, $P$, $K$, $J$) con unidades físicas documentadas.
\item \texttt{results/} regenerable con semilla fija.
\end{enumerate}

\section{Referencia rápida de fórmulas}

\begin{align}
\Dq(r) &= \Dq_0\,(r_0/r)^5\\
E_{\HS} &= -4\,\Dq(r)\\
E_{\LS} &= -24\,\Dq(r) + P\\
E_{\mathrm{elast}} &= \frac{K}{2}\sum_{(i,j)}(r_{ij}-r_i-r_j)^2\\
E_{\mathrm{total}} &= \sum_i E_{\mathrm{el}} + E_{\mathrm{elast}} \quad \bigl(+\text{opcional } -J\sum\sigma_i\sigma_j\bigr)\\
p_{\mathrm{aceptación}} &= \min\!\bigl(1,\ e^{-\Delta E/\kb T}\bigr)
\end{align}

Bibliografía: \texttt{referencias/README.md}.

\end{document}
